\documentclass[11pt,a4paper]{article}

\usepackage[utf8]{inputenc}
\usepackage[T1]{fontenc}
\usepackage{lmodern}
\usepackage[english]{babel}
\usepackage{amsmath,amssymb,amsthm}
\usepackage[expansion=false]{microtype}
\usepackage{booktabs}
\usepackage{siunitx}
\usepackage{algorithm}
\usepackage{algpseudocode}
\usepackage{geometry}
\geometry{margin=2.5cm}
\usepackage{hyperref}

\newtheorem{theorem}{Theorem}
\newtheorem{proposition}[theorem]{Proposition}
\theoremstyle{definition}
\newtheorem{definition}{Definition}

\hypersetup{
  colorlinks=true,
  linkcolor=black,
  citecolor=blue,
  urlcolor=blue,
  pdftitle={A computer-verified lower bound for OEIS A338700(9)},
  pdfauthor={Carlo Corti}
}

\sisetup{group-separator={\,},group-minimum-digits=4}

\title{A computer-verified lower bound for the ninth term of OEIS sequence A338700}

\author{Carlo Corti \\[0.4ex]
\small Independent researcher \\[0.2ex]
\small \href{mailto:carlo.corti@outlook.com}{carlo.corti@outlook.com}}

\date{May 4, 2026}

\begin{document}

\maketitle

\begin{abstract}
The OEIS sequence A338700 records, for each $n\ge 1$, the smallest prime
$p_k$ such that $n$ consecutive primes $p_k, p_{k+1}, \ldots, p_{k+n-1}$
are all Sophie Germain primes. Eight terms are known, the largest being
$a(8) = 1\,372\,604\,395\,439$ (Ehrenstein, 2021). We report an exhaustive
computer search of the interval $[a(8),\,5\cdot 10^{14}]$ that did not
yield any new term. Consequently we establish the verified lower bound
\[
  a(9) \;>\; 5 \cdot 10^{14}.
\]
The search processed $\num{15186821470495}$ candidate primes over a total
wall-clock time of 52~h 49~min on a 16-thread consumer-class machine
(AMD Ryzen 9~7940HS). The source code, selected campaign log, and reproducibility artefacts are released under the MIT licence~\cite{repo}.
\end{abstract}

\section{Definition and known terms}

Let $p_1 < p_2 < p_3 < \ldots$ denote the primes in increasing order.
A prime $p$ is called a \emph{Sophie Germain prime} if $2p+1$ is also
prime; the set of such primes is OEIS A005384~\cite{oeis-a005384}.
Long runs of consecutive Sophie Germain primes are rare and their
distribution is governed by a Hardy--Littlewood type heuristic~\cite{hardy-littlewood}
(§\ref{sec:campaign}). Each new term of A338700 represents a
record-length run, and exact computation of $a(9)$ would be the first
extension of the sequence since Ehrenstein's 2021 contribution of $a(8)$.
The present work does not find $a(9)$ but establishes a substantially
improved lower bound.

\begin{definition}[OEIS A338700~\cite{oeis-a338700}]
For $n \ge 1$, $a(n)$ is the smallest prime $p_k$ such that the $n$
consecutive primes $p_k, p_{k+1}, \ldots, p_{k+n-1}$ are all Sophie
Germain primes.
\end{definition}

The function $n \mapsto a(n)$ is non-decreasing: any run of length $n+1$
contains a run of length $n$ at the same starting prime, hence
$a(n+1) \ge a(n)$, with equality possible if the first run of length $n$
already happens to extend to length $n+1$.

The eight known terms are listed in Table~\ref{tab:known}.

\begin{table}[ht]
\centering
\caption{Known terms of A338700.}
\label{tab:known}
\begin{tabular}{rrl}
\toprule
$n$ & $a(n)$ & Source \\
\midrule
1 & 2 & trivial \\
2 & 2 & trivial \\
3 & 2 & trivial \\
4 & \num{1433849} & OEIS \\
5 & \num{9816899} & OEIS \\
6 & \num{445480319} & OEIS \\
7 & \num{298098924131} & G.\ Resta, 2004~\cite{rivera-puzzle71} \\
8 & \num{1372604395439} & M.\ Ehrenstein, 2021~\cite{oeis-a338700} \\
\bottomrule
\end{tabular}
\end{table}

\section{Search strategy}

\subsection{Reduction to a sieve problem}

Finding $a(n)$ for a given $n$ amounts to scanning the primes in
increasing order and maintaining the length of the current run of
consecutive Sophie Germain primes; the first index at which the run
reaches length $n$ yields $a(n)$. Two equivalent computational kernels
are available:

\begin{itemize}
\item \textbf{MR mode.} Generate the primes $p$ in $[L,U]$ with a
segmented sieve of Eratosthenes, then test each $p$ for the Sophie
Germain property by a Miller--Rabin primality test~\cite{miller,miller-rabin} on $2p+1$.

\item \textbf{Double-sieve mode.} Sieve simultaneously the interval
$p \in [L,U]$ (yielding a bitmap $P$ of primes) and the interval
$q \in [2L+1,\,2U+1]$ (yielding a bitmap $Q$ of primes). The Sophie
Germain primes in $[L,U]$ are then exactly the indices for which both
$P[p]$ and $Q[2p+1]$ are set. No Miller--Rabin call is needed.
\end{itemize}

For the range of interest the double-sieve mode is faster because
Miller--Rabin on a $64$-bit operand costs several hundred CPU cycles,
whereas a sieve strike is a single bit operation. The production
campaign used double-sieve mode throughout.

Both bitmaps are sieved using base primes up to
$\lceil\sqrt{2U+1}\,\rceil$. For the production campaign
($U = 5\cdot 10^{14}$) this bound is approximately $3.16\cdot 10^7$;
the resulting list fits in roughly 14~MB and is shared read-only across
all threads.

\subsection{Numerical types}
\label{sec:numerical-types}

For $p < 2^{63}$ one has $2p+1 < 2^{64}$, so all prime arithmetic fits in
\texttt{uint64\_t}. Modular multiplications inside Miller--Rabin
(retained for the validation harness) use \texttt{\_\_uint128\_t}
intermediate products, which GCC compiles to native \texttt{MUL/DIV}
instructions on x86-64 without library calls.

The Miller--Rabin implementation is \emph{deterministic} for all inputs
below $3.317\cdot 10^{24}$ using the twelve witness bases
$\{2, 3, 5, 7, 11, 13, 17, 19, 23, 29, 31, 37\}$, which covers the
full range $n < 2^{64}$ relevant to this search. This
deterministic bound is established in~\cite{mr-deterministic}.

\subsection{Algorithm and block-boundary reconciliation}
\label{sec:algorithm}

The interval $[L,U]$ is partitioned into contiguous \emph{blocks} of fixed
width $W = 2^{21}$ odd integers (i.e.\ each block $[\ell_i, h_i]$ spans
$W$ consecutive odd integers). Within each block, the per-block kernel
(Algorithm~1) sieves both $P$ and $Q$, walks the primes in $[\ell_i, h_i]$
in increasing order, and emits a fixed-size summary $s_i$ together with all
runs entirely contained in the block. After the parallel phase
completes, a deterministic serial pass (Algorithm~2) walks the
summaries $s_0, s_1, \ldots$ in order and reconstructs the runs that
cross block boundaries.

Two integer parameters govern run tracking throughout:
\begin{itemize}
\item $n_{\min}$ (\texttt{cfg\_target}): the target run length, equal
to~$9$ for the production campaign. A run is recorded only when its
length reaches or exceeds $n_{\min}$.
\item $M$ (\texttt{max\_run}): the run-length cap, set equal to
$n_{\min}$ in the production campaign. Runs longer than $M$ are
truncated to $M$ in the carry and prefix/suffix fields; since any run
of length $\ge n_{\min}$ disproves the lower bound, this truncation
does not affect the correctness of the search.
\end{itemize}

\paragraph{Per-block summary.} Each block emits a record $s_i$ with
the following fields, where ``Sophie Germain prime'' is abbreviated SG:
\begin{itemize}
\item \texttt{num\_primes}: the number of primes in $[\ell_i, h_i]$;
\item \texttt{first\_prime}, \texttt{last\_prime}: the smallest and
largest primes in $[\ell_i, h_i]$ (or $0$ if the block contains no
primes);
\item \texttt{prefix\_len}, \texttt{prefix\_start}: length and starting
prime of the maximal SG-run that begins at \texttt{first\_prime}, or $0$
if \texttt{first\_prime} is not SG;
\item \texttt{suffix\_len}, \texttt{suffix\_start}: length and starting
prime of the maximal SG-run that ends at \texttt{last\_prime}, or $0$
if \texttt{last\_prime} is not SG;
\item \texttt{all\_sg}: a Boolean flag, true iff every prime in the
block is SG (computed directly in Algorithm~1, line~4 and updated
through the loop).
\end{itemize}

\begin{algorithm}[ht]
\caption{Per-block scan (executed once per block, in parallel)}
\begin{algorithmic}[1]
\Require base primes, block range $[\ell_i, h_i]$, run-length cap~$M$,
recording threshold~$n_{\min}$
\State sieve $P$ on odd integers in $[\ell_i, h_i]$
\State sieve $Q$ on odd integers in $[2\ell_i+1, 2h_i+1]$
\State $\textit{run} \gets 0$;\ \ $\textit{run\_start} \gets 0$
\State zero all fields of $s_i$;\ \ $s_i.\texttt{all\_sg} \gets \textbf{true}$
\ForAll{primes $p$ in $[\ell_i, h_i]$ taken in increasing order}
  \State $s_i.\texttt{num\_primes} \mathrel{+}= 1$
  \If{$s_i.\texttt{first\_prime} = 0$}\ \ $s_i.\texttt{first\_prime} \gets p$\EndIf
  \State $s_i.\texttt{last\_prime} \gets p$
  \State $\textit{sg} \gets P[p] \wedge Q[2p+1]$
  \If{$\textit{sg}$}
    \If{$\textit{run} = 0$}\ \ $\textit{run\_start} \gets p$\EndIf
    \State $\textit{run} \gets \min(\textit{run}+1,\,M)$
    \If{$\textit{run} \ge n_{\min}$}
      \State record run of length $\textit{run}$ at $\textit{run\_start}$ in thread-local record
    \EndIf
  \Else
    \State $s_i.\texttt{all\_sg} \gets \textbf{false}$
    \State $\textit{run} \gets 0$
  \EndIf
  \State maintain $s_i.\texttt{prefix\_*}$ until first non-SG prime; maintain $s_i.\texttt{suffix\_*}$ at every SG-extending step
\EndFor
\end{algorithmic}
\end{algorithm}

\begin{algorithm}[ht]
\caption{Boundary reconciliation (executed serially after the parallel phase)}
\begin{algorithmic}[1]
\Require summaries $s_0, s_1, \ldots, s_{B-1}$ in block order; cap $M$;
recording threshold~$n_{\min}$
\State $\textit{carry} \gets 0$;\ \ $\textit{carry\_start} \gets 0$
\For{$i = 0, 1, \ldots, B-1$}
  \If{$s_i.\texttt{num\_primes} = 0$}\ \ \textbf{continue}\EndIf
  \If{$\textit{carry} > 0 \wedge s_i.\texttt{prefix\_len} > 0$}
     \State $\textit{merged} \gets \min(\textit{carry} + s_i.\texttt{prefix\_len},\,M)$
     \If{$\textit{merged} \ge n_{\min}$}
       \State record run of length $\textit{merged}$ at $\textit{carry\_start}$ in global record
     \EndIf
  \EndIf
  \If{$s_i.\texttt{all\_sg}$}
     \If{$\textit{carry} = 0$}\ $\textit{carry\_start} \gets s_i.\texttt{first\_prime}$\EndIf
     \State $\textit{carry} \gets \min(\textit{carry} + s_i.\texttt{num\_primes},\,M)$
     \If{$\textit{carry} \ge n_{\min}$}
       \State record run of length $\textit{carry}$ at $\textit{carry\_start}$ in global record \Comment{multi-block run}
     \EndIf
  \ElsIf{$s_i.\texttt{suffix\_len} > 0$}
     \State $\textit{carry} \gets s_i.\texttt{suffix\_len}$;\ \ $\textit{carry\_start} \gets s_i.\texttt{suffix\_start}$
  \Else
     \State $\textit{carry} \gets 0$;\ \ $\textit{carry\_start} \gets 0$
  \EndIf
\EndFor
\end{algorithmic}
\end{algorithm}

\paragraph{Invariant and correctness.}
Let $r(p)$ denote the length of the maximal run of consecutive SG primes
ending at the prime $p$, taken over the global prime sequence (not
restricted to any block). The reconciliation algorithm maintains the
following invariant.

\begin{proposition}[Reconciliation invariant]
\label{prop:invariant}
After processing summaries $s_0, \ldots, s_{i-1}$:
\begin{itemize}
\item if the last prime processed so far is SG, then $\textit{carry}$
equals $\min(r(\text{last prime}), M)$ and
$\textit{carry\_start}$ is the starting prime of that run;
\item otherwise $\textit{carry} = 0$.
\end{itemize}
Moreover, every maximal SG-run that ends within $\bigcup_{j<i} [\ell_j, h_j]$
has been recorded in either a thread-local record (Algorithm~1) or the
global record (Algorithm~2), at least once.
(For the purposes of this search, recording a run multiple times is
harmless: the global record retains only the best starting prime per
run length.)
\end{proposition}

\noindent\emph{Empirical verification (sketch).}
The invariant is not formally proved by induction in this manuscript;
correctness is established empirically through the three-level validation
harness (§\ref{sec:validation}) and the overlapping-window cross-tests.
The following sketch explains the design intent. A maximal SG-run either
lies entirely within a single block (recorded by Algorithm~1, line~14) or
crosses one or more block boundaries. In the latter case, its starting
prime lies either at the first prime of some block (captured by
$s_i.\texttt{prefix\_*}$) or at $s_i.\texttt{suffix\_start}$ for the
leftmost block~$i$ in which the run is non-empty. Algorithm~2 closes the
cross-boundary run exactly when its right endpoint occurs (line~5), and
never closes the same run twice because $\textit{carry}$ is reset on the
first non-SG prime encountered afterwards. Multi-block all-SG runs are
tracked by accumulating \textit{carry} across consecutive all-SG blocks
(line~12).

\paragraph{Concurrency discipline.}
The parallel phase has no shared mutable state. Each thread~$t$ writes
exclusively to (i)~its own block summary $s_i$ (the assignment $i \mapsto
t$ is fixed by the OpenMP schedule for the duration of that block), and
(ii)~its own record buffer $\texttt{thread\_rec}[t]$. Both structures
are cache-line aligned (64~bytes) to prevent false sharing. After the
parallel phase, the serial post-pass merges all $\texttt{thread\_rec}[t]$
into a global record and then runs Algorithm~2. There are no locks,
atomics, or barriers within the parallel phase other than the implicit
end-of-region barrier.

\subsection{Validation harness}
\label{sec:validation}

Three validation levels were used to certify the implementation before
the production campaign:

\begin{itemize}
\item \texttt{small}: $[0,\,10^7]$, verifies $a(1),\ldots,a(5)$
(under one second);
\item \texttt{medium}: $[0,\,5\cdot 10^8]$, verifies $a(1),\ldots,a(6)$
(seconds);
\item \texttt{full}: $[0,\,2\cdot 10^{12}]$, verifies $a(1),\ldots,a(8)$
(hours).
\end{itemize}

All three levels reproduce the corresponding OEIS values exactly in both
MR and double-sieve mode; the agreement is guaranteed (not merely
probabilistic) because the MR harness uses the twelve-base deterministic
witness set described in §\ref{sec:numerical-types}. Block-boundary
stress tests (eight overlapping windows, both modes, each tested against
the other) all match byte-for-byte.

The double-sieve operates on odd integers only, so the prime $p = 2$ is
handled as a special case in validation mode: when the search interval
includes $2$ (i.e.\ $L \le 2$), the harness explicitly initialises the
run-tracking state to reflect $p = 2$ as the first Sophie Germain prime
($2p+1 = 5$ is prime) before entering the odd-only sieve at $p = 3$.
This special case applies only to the validation tiers; it is not
exercised by the production campaign, whose lower bound $L = a(8)$ is
odd.

\section{The campaign}
\label{sec:campaign}

\subsection{Search window}

Because A338700 is non-decreasing, $a(9) \ge a(8)$ and the search may
start at $L = a(8) = \num{1372604395439}$. Equality is possible (the run
of length~$8$ starting at $a(8)$ might extend to length~$9$ or more), so
the implementation includes the starting prime $a(8)$ in its scan. The
initial carry is set to~$0$: the run-length analysis is self-contained
within $[L, U]$ and does not depend on primes below~$L$.

The upper bound $U = 5\cdot 10^{14}$ was chosen as a round value
approximately $360$ times larger than $a(8)$, reachable within a
reasonable wall-clock budget on the available hardware. Under the
Hardy--Littlewood heuristic~\cite{hardy-littlewood} the median prediction
for $a(9)$ is of order $10^{14}$, so $U$ exceeds it by a factor
of~$\sim 5$. With $L \approx 1.37\cdot 10^{12}$ this gives a
search interval of width $\sim 4.99\cdot 10^{14}$, of order $10^{14}$.

\subsection{Hardware and software}

\begin{itemize}
\item CPU: AMD Ryzen 9~7940HS (Zen 4, 8 cores / 16 threads).
\item RAM: 64\,GB DDR5.
\item OS: Linux Mint 22.3.
\item Compiler: GCC 13.3.0, C17, flags
\verb|-O2 -march=native -fopenmp -std=c17|.
\item OpenMP schedule: \texttt{OMP\_SCHEDULE=guided,4}.
\item Implementation: 16 threads, 16-way parallel double-sieve, segment
of 2\,097\,152 odd candidates per block, batches of 4\,096 blocks each.
\item Licence: MIT.
\end{itemize}

The full source code, build system, validation harness, and selected
campaign log are released in the repository~\cite{repo}.

\subsection{Results}

The campaign was carried out in ten consecutive sessions, with periodic
checkpoints to a state file allowing exact resume. Aggregate figures:

\begin{table}[ht]
\centering
\caption{Campaign aggregates.}
\label{tab:campaign}
\begin{tabular}{lr}
\toprule
Search window & $[\num{1372604395439},\,5\cdot 10^{14}]$ \\
Total wall-clock time (10 sessions) & 52~h 49~min 11~s \\
Total batches processed & \num{29024} \\
Primes scanned & \num{15186821470495} \\
Peak throughput (warm batch, 16 threads) & ${\sim}1.76\cdot 10^{9}$ odd cand./s \\
Average throughput (full campaign) & ${\sim}8.0\cdot 10^{7}$ primes/s \\
Sophie Germain primes found in window & \num{598620723232} \\
Longest run found in window & length 8, at $p = a(8)$ \\
New term $a(9)$ & not found \\
\bottomrule
\end{tabular}
\end{table}

The peak figure refers to the throughput of odd-integer candidates processed
by the double-sieve in the warmest batches of the final session; it is
measured in sieve slots per second, not confirmed primes per second. The
average throughput is computed from total primes scanned divided by total
wall-clock time (190\,151\,s) and reflects the overhead of ten session
starts, periodic checkpoint writes during execution, and the natural
variation in prime density across the interval.

The proportion of Sophie Germain primes among the primes in the window,
\[
  \frac{\num{598620723232}}{\num{15186821470495}}
  \;\approx\; 0.0394,
\]
is consistent with the Hardy--Littlewood prediction~\cite{hardy-littlewood}
$\pi_{\mathrm{SG}}(x)/\pi(x) \sim 2 C_2 / \log x$ with twin-prime
constant $C_2 \approx 0.6601618$ at the relevant scale
$\log x \approx 33$, which gives the heuristic value $\approx 0.0400$.

\section{Result}

The exhaustive search produced the following verified statement.

\begin{theorem}
Let $a(n)$ denote the $n$-th term of OEIS A338700. Then
\[
  a(9) \;>\; 5 \cdot 10^{14}.
\]
\end{theorem}

\begin{proof}[Verification]
Every prime in the closed interval $[a(8),\,5\cdot 10^{14}]$ was
classified as Sophie Germain or not, by means of the double-sieve
procedure described in §2. The longest run of consecutive Sophie
Germain primes recorded in this interval has length~$8$ and starts at
$p = a(8)$; no run of length~$9$ exists in the interval. The
implementation reproduces $a(1),\ldots,a(8)$ exactly under three
independent validation levels including a full re-derivation up to
$2\cdot 10^{12}$. Block-boundary reconciliation is verified by
overlapping-window cross-tests and by equivalence between two
independent kernels (MR and double-sieve). The selected campaign log and source code are archived at~\cite{repo}. The lower bound
has been communicated to the OEIS and is recorded at~\cite{corti-oeis}.
\end{proof}

\section{Reproducibility}

The repository~\cite{repo} contains:

\begin{itemize}
\item the full C17 source (single binary, $<3000$~LOC);
\item a \texttt{Makefile} with five main targets:
\texttt{validate-small}, \texttt{validate-medium},
\texttt{validate-full}, \texttt{bench-schedule},
\texttt{release-znver4};
\item the selected campaign log
(\texttt{a338700\_campaign.log}), documenting the main campaign sessions.
\end{itemize}

The exact compilation flags, OpenMP schedule, and block parameters used
in the production binary are recorded in the repository README.

\subsection{Edge-block handling}
The search interval $[L,U]$ is padded internally to a multiple of
$\texttt{block\_odds} = 2^{21}$ odd integers. Positions beyond $U$ in the
final block are masked out before the per-block scan emits its summary,
so that all summary fields refer only to primes in $[L,U]$. The sieve
itself does no bounds-checking on the padded region; correctness rests
on the masking step.

\subsection{Checkpoint format and resume integrity}
The checkpoint file holds a single $\texttt{state\_t}$ structure with
the following fields, in order:
\begin{itemize}
\setlength{\itemsep}{0pt}
\item a 32-bit magic number and a 32-bit format version (currently~$2$);
\item the search configuration:
{\sloppy $\texttt{cfg\_low}$, $\texttt{cfg\_high}$,
$\texttt{cfg\_block\_odds}$, $\texttt{cfg\_batch\_blocks}$,
$\texttt{cfg\_target}$, $\texttt{cfg\_max\_run}$, $\texttt{cfg\_mr}$;\par}
\item the next-batch index;
\item the carry pair $(\textit{carry}, \textit{carry\_start})$ from
Proposition~\ref{prop:invariant};
\item the running totals $\texttt{total\_primes}$ and $\texttt{total\_sg}$;
\item the global record $\texttt{global\_rec}$ (best run-starting primes
per length);
\item a 64-bit checksum.
\end{itemize}

The checksum is a byte-wise 8-lane XOR over the structure up to (but
excluding) the checksum field itself. On resume, the file is rejected
if the magic, version, or checksum does not match. Configuration fields
loaded from the file overwrite the runtime configuration, so the resumed
run uses exactly the same search parameters as the interrupted one.
Writes are atomic: the new state is written to a temporary file, fsynced,
and renamed over the previous state file, after which the parent
directory is fsynced as well.

\section*{AI use disclosure}

This work was carried out with substantial support from large language
models, used as a research assistant throughout the full workflow.
Specifically: mathematical analysis of the candidate space and
Hardy--Littlewood heuristics (Claude Opus 4.7, Anthropic); design of
the computational strategy (Claude Opus 4.7); code generation in C17
(GPT-5.5, OpenAI); LLM-based code review of successive versions, with
six independent model passes (Claude, ChatGPT, DeepSeek, Gemini, Grok,
Qwen); and composition of this manuscript (Claude Sonnet 4.6,
Anthropic). Every quantitative claim in this paper was verified
computationally by the author; all mathematical statements were checked
by the author before inclusion. The author takes full scientific
responsibility for the content. The exact model versions and API
identifiers used at each stage are recorded in the repository
at~\cite{repo}.

\begin{thebibliography}{9}

\bibitem{oeis-a338700}
The OEIS Foundation Inc.,
\emph{Sequence A338700: Smallest prime $p$ such that $p$ and the $n-1$
following primes are all Sophie Germain primes},
The On-Line Encyclopedia of Integer Sequences,
\url{https://oeis.org/A338700}.
Sequence created by M.\ Sariyar, April 24, 2021; term $a(8)$
contributed by M.\ Ehrenstein, April 29, 2021.

\bibitem{oeis-a005384}
The OEIS Foundation Inc.,
\emph{Sequence A005384: Sophie Germain primes $p$: $2p+1$ is also prime},
The On-Line Encyclopedia of Integer Sequences,
\url{https://oeis.org/A005384}.

\bibitem{rivera-puzzle71}
C.\ Rivera,
\emph{Puzzle 71: Consecutive primes and Cunningham chains},
The Prime Puzzles \& Problems Connection,
\url{https://www.primepuzzles.net/puzzles/puzz_071.htm}.

\bibitem{hardy-littlewood}
G.~H.\ Hardy and J.~E.\ Littlewood,
\emph{Some problems of `Partitio Numerorum'\,III: On the expression of a
number as a sum of primes},
Acta Mathematica \textbf{44} (1923), 1–70.
DOI: \href{https://doi.org/10.1007/BF02403921}{10.1007/BF02403921}.

\bibitem{miller}
G.~L.\ Miller,
\emph{Riemann's hypothesis and tests for primality},
Journal of Computer and System Sciences \textbf{13} (1976), 300–317.
DOI: \href{https://doi.org/10.1016/S0022-0000(76)80043-8}{10.1016/S0022-0000(76)80043-8}.

\bibitem{miller-rabin}
M.~O.\ Rabin,
\emph{Probabilistic algorithm for testing primality},
Journal of Number Theory \textbf{12} (1980), 128–138.
DOI: \href{https://doi.org/10.1016/0022-314X(80)90084-0}{10.1016/0022-314X(80)90084-0}.

\bibitem{mr-deterministic}
J.~Sorenson and J.~Webster,
\emph{Strong pseudoprimes to twelve prime bases},
Mathematics of Computation \textbf{86} (2017), 985–1003.
DOI: \href{https://doi.org/10.1090/mcom/3134}{10.1090/mcom/3134}.
The twelve bases $\{2,3,5,7,11,13,17,19,23,29,31,37\}$ yield a
deterministic test for all $n < 3.317\cdot 10^{24}$, covering the
full \texttt{uint64\_t} range used in this work.

\bibitem{corti-oeis}
C.\ Corti,
\emph{Contribution to OEIS A338700: lower bound $a(9) > 5\cdot 10^{14}$},
The On-Line Encyclopedia of Integer Sequences,
\url{https://oeis.org/A338700}, 2026.

\bibitem{repo}
C.\ Corti,
\emph{A338700 — extending the OEIS sequence: source code, logs and
reproducibility artefacts},
GitHub repository: \url{https://github.com/carcorti/A338700}.
Archived release: Zenodo, DOI \href{https://doi.org/10.5281/zenodo.20039591}{\texttt{10.5281/zenodo.20039591}}.

\end{thebibliography}

\end{document}
